\chapter{The model}

\begin{definition}[Block covariance matrix, \S2.1]\label{def:SB}
  \lean{RBM.SB}\leanok
  Let $L\ge 3$ and work on $\ZL$.  Set
  \[ \SB_{ab} \;=\; \tfrac13\,\mathbf 1\!\left(\operatorname{dist}_{\ZL}(a,b)\le 1\right),
     \qquad a,b\in\ZL . \]
  Equivalently $\SB$ is the circulant matrix generated by the kernel supported on
  $\{0,1,-1\}$ with value $1/3$.
\end{definition}

\begin{remark}\label{rem:L3}
  The paper leaves the hypothesis $L\ge 3$ implicit.  It is needed: for $L=1$ the
  row sums are $1/3$, and for $L=2$ one has $1=-1$ in $\ZL$ so the row sums are
  $2/3$.  See \texttt{docs/paper-deltas.md}.
\end{remark}

\begin{lemma}[Elementary properties of $\SB$]\label{lem:SB-basic}
  \lean{RBM.SB_isSymm, RBM.SB_apply_add_right, RBM.sum_SB_row, RBM.SB_mulVec_one}
  \leanok\uses{def:SB}
  $\SB$ is symmetric and translation invariant, each row sums to $1$, and
  $\SB \mathbf 1 = \mathbf 1$.
\end{lemma}
\begin{proof}\leanok\uses{def:SB}
  Symmetry and translation invariance are immediate from the circulant form.  The
  row sums reduce, after the change of variables $b\mapsto a-b$, to the total mass
  of the kernel, which is $3\cdot\frac13=1$ because $0,1,-1$ are distinct in $\ZL$
  when $L\ge 3$.
\end{proof}

\chapter{The propagator}

\begin{definition}[Propagator, Definition 2.13]\label{def:Theta}
  \lean{RBM.Theta}\leanok\uses{def:SB}
  For $\xi\in\CC$ with $|\xi|<1$ set $\Th_\xi := (1-\xi\SB)^{-1}$.
\end{definition}

\begin{lemma}[Operator norm of $\SB$]\label{lem:norm-SB}
  \lean{RBM.norm_SB}\leanok\uses{lem:SB-basic}
  $\|\SB\|_{\ell^\infty\to\ell^\infty}=1$.
\end{lemma}
\begin{proof}\leanok\uses{lem:SB-basic}
  The $\ell^\infty$ operator norm is the maximum $\ell^1$ row norm; all entries are
  non-negative and every row sums to $1$.
\end{proof}

\begin{lemma}[$\Th_\xi$ is a two-sided inverse]\label{lem:Theta-inv}
  \lean{RBM.Theta_mul, RBM.mul_Theta, RBM.eq_Theta_of_mul}\leanok
  \uses{def:Theta, lem:norm-SB}
  For $|\xi|<1$ the matrix $1-\xi\SB$ is invertible and $\Th_\xi$ is its two-sided
  inverse.  Moreover any left inverse of $1-\xi\SB$ equals $\Th_\xi$.
\end{lemma}
\begin{proof}\leanok\uses{lem:norm-SB}
  $\|\xi\SB\|=|\xi|<1$, so the geometric series converges and $1-\xi\SB$ is a unit.
\end{proof}

\begin{lemma}[Lemma 2.14 (1): symmetry]\label{lem:Theta-sym}
  \lean{RBM.Theta_transpose}\leanok\uses{lem:Theta-inv, lem:SB-basic}
  $(\Th_\xi)_{xy}=(\Th_\xi)_{yx}$.
\end{lemma}
\begin{proof}\leanok\uses{lem:Theta-inv}
  $\Th^{\mathsf T}(1-\xi\SB) = \Th^{\mathsf T}(1-\xi\SB)^{\mathsf T}
   = \bigl((1-\xi\SB)\Th\bigr)^{\mathsf T} = 1$, so $\Th^{\mathsf T}$ is a left
  inverse; conclude by uniqueness.
\end{proof}

\begin{lemma}[Lemma 2.14 (2): translation invariance]\label{lem:Theta-transl}
  \lean{RBM.Theta_apply_add_right}\leanok\uses{lem:Theta-inv, lem:SB-basic}
  $(\Th_\xi)_{xy}=(\Th_\xi)_{x+1,y+1}$.
\end{lemma}
\begin{proof}\leanok\uses{lem:Theta-inv}
  Conjugating by the shift, i.e.\ taking the submatrix along $a\mapsto a+c$,
  fixes $1-\xi\SB$; conclude by uniqueness of the inverse.
\end{proof}

\begin{lemma}[Lemma 2.14 (3): commutativity]\label{lem:Theta-comm}
  \lean{RBM.Theta_commute_SB, RBM.Theta_commute}\leanok\uses{lem:Theta-inv}
  $[\SB,\Th_\xi]=[\Th_{\xi'},\Th_\xi]=0$.
\end{lemma}
\begin{proof}\leanok\uses{lem:Theta-inv}
  Inverses of commuting elements commute.
\end{proof}

\begin{lemma}[Lemma 2.14 (5): random walk representation]\label{lem:Theta-series}
  \lean{RBM.Theta_eq_tsum}\leanok\uses{lem:Theta-inv}
  $\Th_\xi=\sum_{k\ge 0}\xi^k(\SB)^k$, convergent in operator norm.
\end{lemma}
\begin{proof}\leanok\uses{lem:Theta-inv}
  Neumann series.
\end{proof}

\begin{lemma}[Row sums]\label{lem:Theta-rowsum}
  \lean{RBM.sum_Theta_row, RBM.Theta_mulVec_one}\leanok
  \uses{lem:Theta-inv, lem:SB-basic}
  $\sum_b (\Th_\xi)_{ab} = (1-\xi)^{-1}$.  Used in (3.47)--(3.48) and in Step 4 of
  the proof of Lemma 3.6.
\end{lemma}
\begin{proof}\leanok\uses{lem:SB-basic}
  Apply $\Th_\xi$ to the constant vector and use $\SB\mathbf 1=\mathbf 1$.
\end{proof}

\begin{lemma}[Equation (2.51)]\label{lem:Theta-deriv}
  \lean{RBM.hasDerivAt_Theta_apply, RBM.Theta_sub_Theta}\leanok
  \uses{lem:Theta-inv}
  $\partial_\xi \Th_\xi = \Th_\xi \SB \Th_\xi$, entrywise.
\end{lemma}
\begin{proof}\leanok\uses{lem:Theta-inv}
  The resolvent identity
  $\Th_\zeta-\Th_\xi=(\zeta-\xi)\,\Th_\zeta\SB\Th_\xi$
  together with continuity of $\xi\mapsto\Th_\xi$.
\end{proof}

\begin{lemma}[$\ell^1$ and $\ell^\infty$ bounds, (3.35)--(3.36)]\label{lem:Theta-bounds}
  \lean{RBM.norm_Theta_le, RBM.sum_norm_Theta_row_le, RBM.norm_Theta_apply_le}\leanok
  \uses{lem:Theta-series, lem:norm-SB}
  $\max_a\sum_b |(\Th_\xi)_{ab}| \le (1-|\xi|)^{-1}$, and in particular
  $|(\Th_\xi)_{ab}| \le (1-|\xi|)^{-1}$.
\end{lemma}
\begin{proof}\leanok\uses{lem:Theta-series}
  Sum the Neumann series in operator norm.
\end{proof}

\begin{lemma}[Band structure and a first decay bound]\label{lem:Theta-band}
  \lean{RBM.SB_pow_apply_eq_zero, RBM.Theta_eq_geom_add, RBM.norm_Theta_apply_le_pow}\leanok
  \uses{lem:Theta-series, lem:Theta-bounds}
  $(\SB)^k$ vanishes outside the band of width $k$, and consequently
  \[ \bigl|(\Th_\xi)_{xy}\bigr| \;\le\;
     \frac{|\xi|^{\operatorname{dist}(x,y)}}{1-|\xi|} . \]
  This is exponential decay with length $1/\log(1/|\xi|)$, weaker than (2.52).
\end{lemma}
\begin{proof}\leanok\uses{lem:Theta-series}
  Truncate the Neumann series at order $d=\operatorname{dist}(x,y)$:
  $\Th_\xi=\sum_{k<d}(\xi\SB)^k+(\xi\SB)^d\Th_\xi$, and the finite sum has
  vanishing $(x,y)$ entry.
\end{proof}

\begin{lemma}[Lemma 2.14 (4): exponential decay, (2.52)]\label{lem:Theta-decay}
  \uses{lem:Theta-band}
  With $\hat\ell(\xi):=\min(|1-\xi|^{-1/2},L)$,
  \[ \bigl|(\Th_\xi)_{xy}\bigr| \;\le\;
     \frac{C\,e^{-c|x-y|/\hat\ell(\xi)}}{|1-\xi|\cdot\hat\ell(\xi)} . \]
  \emph{Not yet formalized.}  The plan is to bypass the contour-shift and Poisson
  summation argument of Appendix B: for the nearest-neighbour $\SB$ in $d=1$ the
  inverse has a closed form.  Writing $\rho=\rho(\xi)$ for the root of
  $\rho^2-(3/\xi-1)\rho+1=0$ of modulus $<1$, one has
  $(\Th_\xi)_{xy}=c(\xi)\bigl(\rho^{\,d}+\rho^{\,L-d}\bigr)$ with $d$ the
  representative of $x-y$ in $[0,L)$, and (2.52)--(2.54) follow by elementary
  estimates on $\rho^d$.
\end{lemma}

\begin{lemma}[Difference estimates (2.53)--(2.54)]\label{lem:Theta-diff}
  \uses{lem:Theta-decay}
  \emph{Not yet formalized.}
\end{lemma}

\chapter{The primitive equation and $\Kcal$}

\begin{definition}[Loops and cut-and-glue operators, Definitions 2.9--2.10]\label{def:loops}
  \emph{Not yet formalized.}  The operators
  $\Gcal^{(a)}_k$, $\Gcal^{(a),L}_{k,l}$, $\Gcal^{(b),R}_{k,l}$ act on the index
  data $(\boldsymbol\sigma,\mathbf a)$.
\end{definition}

\begin{definition}[Primitive equation, Definition 2.12]\label{def:primitive}
  \uses{def:loops, def:Theta}
  \emph{Not yet formalized.}
\end{definition}

\begin{lemma}[Classification of canonical partitions, Lemma 3.2]\label{lem:3.2}
  \uses{def:loops}
  \emph{Not yet formalized.}  Canonical partitions of the oriented $n$-gon are in
  bijection with crossing-free subsets of $\mathbb{Z}_n^{\mathrm{off}}$.
\end{lemma}

\begin{lemma}[Tree representation of $\Kcal$, Lemma 3.4]\label{lem:3.4}
  \uses{lem:3.2, def:primitive, lem:Theta-deriv}
  \emph{Not yet formalized.}  Formalization strategy: take the tree sum as the
  \emph{definition}, prove it satisfies (2.48), and invoke Gr\"onwall uniqueness,
  rather than developing existence theory for the ODE.
\end{lemma}

\begin{lemma}[Ward identity for $\Kcal$, Lemma 3.6]\label{lem:3.6}
  \uses{lem:3.4, lem:Theta-rowsum}
  \emph{Not yet formalized.}  The difference $D$ of the two sides satisfies a
  homogeneous linear ODE with zero initial value.
\end{lemma}

\begin{lemma}[Sum-zero property, Lemma 3.10]\label{lem:3.10}
  \uses{lem:3.6, lem:Theta-decay}
  \emph{Not yet formalized.}
\end{lemma}

\begin{lemma}[Bound on $\Kcal$, Lemma 3.11]\label{lem:3.11}
  \uses{lem:3.10}
  \emph{Not yet formalized.}
\end{lemma}

\chapter{The stochastic layer (out of scope for now)}

\begin{remark}\label{rem:stochastic}
  Sections 2.4, 2.6, 2.7 and 5--7 of the paper rest on It\^o calculus for matrix
  Brownian motion, which Mathlib does not have: there is no It\^o formula, no
  matrix-valued SDE theory, and no Dyson Brownian motion.  The loop hierarchy
  (2.45), Lemmas 2.18--2.20 and Theorem 2.21 are therefore left as stated
  hypotheses, so that the logical skeleton of the main theorems still typechecks.
\end{remark}
